\begin{proposition}[the cap Pontryagin pairing]\label{push:src:C.3.19}
For the actual cap $M_\psi$ and meridian class $x$ of \ref{push:src:C.3.14},
\[
 \bigl|\langle p_1(M_\psi)x,[M_\psi]\rangle\bigr|=24.
\]
Thus the framing coordinate $\beta'$ in Wall's convention has
$|\beta'|=6$. The orientation sign and the cubic
$k_\psi=\langle x^3,[M_\psi]\rangle$ are not determined by
this proposition, and no numerical $J_\psi$ or monodromy class
is asserted.
\end{proposition}
\Uses{the hypotheses, notation and definitions stated here; \cref{push:src:C.3.14}; \cref{push:src:C.3.18}}

\begin{proof}
We use the cohomology suspension of the actual universal class
$p_1$, denoted $s(p_1)\in H^3(SO;\Z)$, with the clutching
convention fixed once throughout. In particular we do not
replace it by an unspecified primitive generator of
$H^3(SO;\Z)$ or of $H^3(\operatorname{Spin};\Z)$.
All computations below may first be made in real cohomology;
they determine the integral characteristic number, and torsion
cannot contribute to its pairing.

Write $N=S^2\times D^4$ and $Y=\partial N=S^2\times S^3$,
so that the cap is $N\cup_{\widetilde\psi}(-N)$. Choose a
stable tangent frame on $N$. Its restriction supplies a stable
frame on the outer copy $P=S^2\times D^3\subset Y$. The
stable frame on $Q=S^5$, restricted to $P$, is homotopic to
this one: the difference is a map $P\to SO$, and
$P\simeq S^2$ while $\pi_2(SO)=0$. Thus the frame comparison
period obtained in \ref{push:src:C.3.18} is also the one computed from the
chosen frame on $N$.

Let $\mathcal A=L\times I\subset P$ be the annular support,
where
\[
 L=S^1_\gamma\times\operatorname{Nil}^3(u,w,q).
\]
Use a global affine nilmanifold frame on $\mathcal A$.
The nontrivial $q$-monodromy in $u,w$ is accounted for by
the invariant nilmanifold frame; on any fixed $q$-angle
three-torus it differs from the marked period frame only
by a constant triangular change. Translation in $\gamma$
preserves this global frame. Let
$A:\mathcal A\to SO$ be the stable change-of-frame map
from this frame to the restricted global frame.

In the affine frame, the derivative of the annular map
$(r,p)\mapsto(r,\rho_{\chi(r)}p)$ is the identity plus
the single shear $\chi'(r)\,dr\otimes\partial_\gamma$.
Multiplying the shear by a parameter from one to zero
contracts it to the identity through invertible matrices,
fixed on both end collars. Therefore the actual derivative
clutching map is homotopic, relative to the complement of
$\mathcal A$, to
\[
 G(r,p)=A(r,\rho_{\chi(r)}p)^{-1}A(r,p).
\]
At both ends $G=1$, since the translation is respectively
zero or a full period. Replacing $A(r,p)$ by its restriction
to one middle slice gives a homotopy of $G$ relative to
these ends. We may consequently write $A=A(p)$.

The Mayer--Vietoris clutching formula gives
\[
 p_1(M_\psi)=\delta\bigl(G^*s(p_1)\bigr),
\]
where $G$ is extended by the identity on the rest of $Y$.
With the boundary orientation associated to this connecting
map, its cup-product pairing is exactly
\[
 \langle p_1(M_\psi)x,[M_\psi]\rangle
 =\langle G^*s(p_1)\smile x|_Y,[Y]\rangle.
\]
Equivalently this is the relative pairing on
$(\mathcal A,\partial\mathcal A)$. Changing a global
clutching orientation reverses both associated signs;
we will retain the absolute value in the numerical conclusion.

We calculate this relative class explicitly. Define
\[
 \Gamma:S^1\times L\longrightarrow SO,
 \qquad\Gamma(t,p)=A(\rho_t p)^{-1}A(p).
\]
It is the identity on $\{0\}\times L$. The suspension
$s(p_1)$ is primitive under multiplication. Slanting its
pullback by the oriented $t$-circle therefore gives
\[
 \int_{S^1}\Gamma^*s(p_1)
 =-\int_{S^1}\rho^*\bigl(A^*s(p_1)\bigr)
 \quad\hbox{in }H^2(L;\R).
\]
There is no additional relative lift ambiguity here:
restriction $H^2(S^1\times L)\to H^2(\{0\}\times L)$
is surjective. Thus $H^3(S^1\times L,\{0\}\times L)$
is exactly the degree-two slant summand. Pulling back by
$t=\chi(r)$ identifies the relative annular class with
the suspension of the right-hand side. This is the
relative trace calculation, including the full-period
endpoint; an absolute equality on $\mathcal A$ alone
would not suffice.

Put $g=\theta_\gamma$ and $a=\arg q/(2\pi)$. The meridian
Euler class has, up to the already fixed overall circle
orientation, the representative
\[
 e=x|_L=dg\wedge da.
\]
Indeed the positive $q$ period has
$A_0\gamma=\gamma-\delta$ and $A_0u=u+w$.
Using coefficient coordinates on the covering mapping
torus, its deck identification can be written
\[
 (g,u,w,d,a)\longmapsto(g,u,w-u,d+g,a+1),
\]
where $d$ is the $\delta$ coordinate. The normalized
connection $\mathrm d d-a\,\mathrm d g$ is invariant and has curvature
$dg\wedge da$. This also shows that no additional
component of the Euler class has been suppressed.

Decompose the real cohomology class
$\alpha=A^*s(p_1)$ using $L=S^1_g\times\operatorname{Nil}^3$:
\[
 \alpha=dg\smile\beta+\alpha_0,
 \qquad\beta\in H^2(\operatorname{Nil}^3;\R),\quad
 \alpha_0\in H^3(\operatorname{Nil}^3;\R).
\]
Translation in $g$ has trace $\beta$. Therefore the preceding
relative calculation and the cup pairing reduce, with a
fixed orientation sign, to
\[
 \pm\int_L dg\wedge da\wedge\beta
 =\pm\int_{T^2_{u,w}}\beta
 =\pm\int_{T^3_{\gamma,u,w}}A^*s(p_1).
\]
The first equality on the right is integration over the
mapping-torus base, with its fibre orientation. Components
of $\beta$ containing $da$, and the component $\alpha_0$,
give zero. This proves that there are no unaccounted
cross terms in the cap pairing.

Finally, the last integral is the relative stable-frame
Pontryagin defect of \ref{push:src:C.3.18}, with at most the inverse
change-of-frame and orientation signs. Its absolute value
is $24$. Hence the cap pairing has absolute value $24$.
Wall's exact formula $\langle p_1x,[M_\psi]\rangle=4\beta'$
from \ref{push:src:C.3.14} gives $|\beta'|=6$, as asserted.
\end{proof}

\begin{remark}[Scope]
The argument uses both the global frame defect furnished by
the bounding four-manifold and the actual Euler class of the
$\delta$ bundle on the annulus. A local unipotent derivative
alone would give neither. The cubic pairing and its orientation
relative to this Pontryagin pairing remain necessary to
evaluate Wall's unframed knot coordinate.
\end{remark}
\Uses{the hypotheses, notation and definitions stated here; the argument and qualifications immediately following}


